% Pursuing Stacks (cotes 134-2 to 134-8): archive cote 134-5, site page 9 (PDF page 10), typescript page 196.
% typescript. Draft compared: ps4.tex:373-441.
% AI-transcribed by Claude Opus 5.5 (claude-opus-5-5, adaptive thinking, effort high), Message Batches API, batch ps2 (msgbatch_01NVwr1Yuo8qJf5rykeGBXPK), 2026-11-12; not yet checked by eye.
\scanpage{134-5/9}
\mspage{196}

29.5.

69. Before resuming \inserted{more} \deleted{with} technical work again, I would like to have a
short retrospective of the last six weeks' scratchwork, now lying on my
desk as a thickly bunch of scratchnotes\hand{,} nobody but I could possibly make
any sense of. \hand{//}The first thing I had on my mind has been there now for
nearly twenty years - ever since it had become clear, in the SGA 5
seminar on L-functions and a propos the formalism of traces in terms of
derived categories, that Verdier's set-up of derived categories was un-%
sufficient for formulating adeq\over{\unsure{ua}}{ua}tely \deleted{even} some rather evident situati-%
ons\unsure{\xout{,}} and relationships, such as the addition formu\hand{\unsure{l}}a for traces, or the
multiplicative formula for determinants. It then became appearant that
the derived category of an abelian category (say) was to\hand{o} coarse an object
in various respects, that it had to be complemented by similar``triangu-%
lated categories''(such as the \xout{\illegible{1}} derived category of a suitable category
of ``triangles'' of complexes, or t\over{\unsure{h}}{h}e whole bunch of derived categories
of categories of filtered complexes of order n with variable n), closely
connected to it. Deligne and Illusie had both set out, independently,
to work out \deleted{a}\inserted{some} set-up meeting the most urgent requirements\hand{\unsure{.}} (Illusie's
treatment in terms of filtered complexes was written down and published
in his thesis six years later (Springer Lecture Notes N\textsuperscript{\underline{o}} 239)). While
adequate for the main tasks then at hand, neither treatment was really
wholly satisfactory to my taste. One main feature I believe making me
feel unconfortable, was that the extra categories which had to be intro-%
duced,to \xout{\unsure{compl}\illegible{1}} \inserted[typed]{round up some}\inserted{what} a stripped-and-naked triangulated categor\over{\unsure{ies}}{y}\hand{,}were
triangulated categories in their own right, in Verdier's sense, but
remaining nearly as stripped \xout{\illegible{1}} by themselves as the initial trian-%
gulated category they were intended to provide clothing for. In other
words, \unsure{theye|there} was a lack of inner stability in the formalism, making it
appear as very much provisional still. Also, while interested in asso-%
ciating to an abelian category a handy sequence of ``filtered derived
categories'', Illusie made no attempt to pin down what exactly the
inner structure of the object he had arrived at was - unlike Verdier,
who \inserted{had} introduced, alongside with the notion of a derived category of an
abelian category, a general notion of triangulated categories, into
which these derived categories would fit. The obvious idea which was in
my head by then for avoiding such shortcomings, was \xout{\illegible{1}}\inserted[typed]{that} an
abelian category $\underline{A}$ gave rise, not only to the \xout{\illegible{1}} single usual
derived category D($\underline{A}$) of Verdier, but also, for every type of diagrams,

\margin{9}
